% [原文 p.133]
设 $G\subset S_k$ 为任意子群。此时，由限制 $\eta\colon R(S_k)\to R(G)$ 从 $\Delta_{n,k}$ 得到的元素
\[
\Delta_{n,k}(G)\in\operatorname{Sym}_k[t_1,\dots,t_n]\otimes R(G)
\]
可以考虑。同样，可以考虑 $\Delta'_{n,k}$ 与
\[
\eta_*\colon R_*(G)\to R_*(S_k)
\]
的合成
\[
\Delta'_{n,k}(G)\colon R_*(G)\to\operatorname{Sym}_k[t_1,\dots,t_n].
\]

特别地，考虑 $k=p$ 为素数、$G=\Z/p$ 为 $p$ 阶循环群的情形。映射
\[
\eta\colon R(S_p)\to R(\Z/p)
\]
的像由 $\Z/p$ 的自明表示 $1$ 与正则表示 $N$ 生成（后者是 $S_p$ 的标准 $p$ 维表示的限制）。因此，由适当的对称函数 $a(t),b(t)$，必须有
\[
\Delta_{n,p}(\Z/p)=a(t)\otimes1+b(t)\otimes N.
\]
在单位元处对 $R(S_p)$ 求值，得到
\[
e_1^p=a+pb.
\]
在 $\Z/p$ 的生成元处求值并利用命题 B.1.7，得到
\[
m_p=a.
\]
因此 $b=\dfrac{e_1^p-m_p}{p}$，并且有
\[
\Bigl(\sum t_i\Bigr)^p\equiv\sum t_i^p \pmod p
\]

% [原文 p.134]
于是 $b$ 的系数当然是整数。这样便证明了下面的命题。

\begin{proposition*}{B.1.10}
  设 $p$ 为素数。此时把 $\Delta_{n,p}$ 从对称群限制到循环群，得到
  \[
  \Delta_{n,p}(\Z/p)=m_p\otimes1+\frac{e_1^p-m_p}{p}\otimes N.
  \]
  其中 $N$ 是 $\Z/p$ 的正则表示。
\end{proposition*}

当 $n\ge p$ 时，通过命题 B.1.2 的同构映射 $\Delta'_n$，把与
\[
\frac{e_1^p-m_p}{p}\in\operatorname{Sym}_p[t_1,\dots,t_n]
\]
对应的元素记作 $\theta^p\in R_*(S_p)$。亦即令
\[
\Delta'_n\theta^p=\frac{e_1^p-m_p}{p}.
\]
此时由命题 B.1.10 可知，$\theta^p$ 是从 $R(S_p)$ 到 $\Z$ 的同态，且限制到 $\Z/p$ 上时给出正则表示 $N$ 的重数。因此，对 $\rho\in R(S_p)$，有
\begin{equation}
\eta(\rho)=\psi^p(\rho)1+\theta^p(\rho)N \tag{B.1.11}
\end{equation}
其中 $\eta\colon R(S_p)\to R(\Z/p)$ 是限制映射。

\section{K 理论中的算子}
\takeaway{本节把 $G$ 向量丛与 $k$ 重张量幂上的对称群作用提升为 $K(X)$ 上的算子：$R_*(S_k)$ 的每个元素给出函子 $K$ 的自然变换，由此得到外幂、对称幂以及 Adams 算子 $\psi^k$。}

设 $X$ 为紧 Hausdorff 空间，$G$ 为有限群。$G$ 作用于 $X$ 时，称 $X$ 为 $G$ 空间。设 $E$ 为 $X$ 上的复向量丛。若 $E$ 是 $G$ 空间且满足下面两个条件，则称 $E$ 为 $G$ 空间 $X$ 上的 $G$ 向量丛。
\begin{enumerate}[label=(\roman*),leftmargin=2.6em,itemsep=1pt]
\item 投影 $E\to X$ 与 $G$ 的作用可交换。
\item 对 $g\in G$，映射 $E_x\to E_{g(x)}$ 是线性的。
\end{enumerate}
$G$ 空间 $X$ 上所有 $G$ 向量丛构成的 Grothendieck 群记作 $K_G(X)$。$X$ 上 $G$ 的作用视为 $X$ 结构的一部分，故在 $K_G(X)$ 中不显式写出。

% [原文 p.135]
对 $G$ 向量丛，可以借 $G$ 上的平均构造不变度量。因此按通常的论证，可以证明短完全序列能与 $G$ 相容地分裂。于是，若
\[
0\longrightarrow E_1\longrightarrow E_2\longrightarrow\cdots\longrightarrow E_n\longrightarrow 0
\]
为 $G$ 向量丛的长完全序列，则 $K_G(X)$ 中欧拉示性数 $\sum(-1)^i[E_i]$ 为零。以这一条为代表的 $K_G(X)$ 的性质，请参看 [4] 与 [9]。

本节只考虑平凡的 $G$ 空间 $X$，即对一切 $x\in X$ 与 $g\in G$ 都有 $g(x)=x$ 的空间。此时，$G$ 向量丛就是给 $X$ 上的向量丛 $E$ 配上同态
\[
G\longrightarrow\Aut E,
\]
其中 $\Aut E$ 是 $E$ 的自同构映射之群。下面考察这样的 $G$ 向量丛。

% [Tutor Note: 这里要求 $X$ 是平凡 $G$ 空间，于是 $g$ 覆盖的底空间映射是恒等，条件 (ii) 自动成为“每个 $g$ 给出 $E$ 的一条纤维自同构”，$G\to\Aut E$ 的表述才有意义。对非平凡 $G$ 空间，$G$ 向量丛只是等变丛，并没有这种同态。]

$G$ 固定的 $E$ 的子空间 $E^G$ 是 $E$ 的子向量丛。实际上，它是投影算子
\[
\frac{1}{|G|}\sum_{g\in G}g
\]
的像，而任意投影算子的像总是子丛 [4]。若 $E$ 与 $F$ 都是 $G$ 向量丛，则 $\Hom(E,F)$ 中与 $G$ 的作用可交换的一切 $\varphi_x\colon E_x\to F_x$ 构成的子空间，是子向量丛 $\Hom_G(E,F)$。事实上，$\Hom_G(E,F)=(\Hom(E,F))^G$。特别地，设 $V$ 为 $G$ 的表示空间，并把对应的 $X$ 上的 $G$ 向量丛 $X\times V$ 也记作 $V$；此时对 $X$ 上任一 $G$ 向量丛 $E$，$\Hom_G(V,E)$ 是向量丛，且有自然的同态映射 $V\otimes\Hom_G(V,E)\to E$。

这里，设 $\{V_\pi\}$ 为 $G$ 的不可约表示的完全集，考虑向量丛的同态
\[
\alpha\colon\sum_{\pi}\bigl\{V_\pi\otimes\Hom_G(V_\pi,E)\bigr\}\longrightarrow E.
\]
对任意 $x\in X$，$\alpha_x$ 是同构。因此，$\alpha$ 是同构映射。

% [原文 p.136]
这样便证明了下面的命题。

\begin{proposition*}{B.2.1}
  若 $X$ 为平凡的 $G$ 空间，则有自然的同构映射
  \[
  K(X)\otimes R(G)\longrightarrow K_G(X).
  \]
\end{proposition*}

特别地，可以把前面的论证应用于向量丛 $E$ 的 $k$ 重张量积 $E^{\otimes k}$ 上 $S_k$ 的自然作用。这样一来，便得到与 $S_k$ 的作用相容的自然分解
\[
E^{\otimes k}\cong\sum_{\pi}\bigl\{V_\pi\otimes\Hom_{S_k}(V_\pi,E^{\otimes k})\bigr\}.
\]

令
\[
\pi(E)=\Hom_{S_k}(V_\pi,E^{\otimes k}),
\]
则 $\pi$ 是向量丛的算子。实际上，$\pi(E)$ 是由分割 $\pi$ 所伴随的 $\mathrm{GL}(n)$（$n=\dim E$）的不可约表示从 $E$ 诱导出的向量丛。然而，此后这一事实并不起特别的作用。

下一步是把这些作用于向量丛的算子扩张为 $K(X)$ 的算子。为此，把 $K(X)$ 表示成集合 $\mathcal{E}(X)$ 按等价关系取商较为方便（$\mathcal{E}(X)$ 的元素起“上链”的作用）。$\mathcal{E}(X)$ 的元素是分次向量丛 $E=\sum_{i\in\Z}E_i$，且除有限个 $i$ 外 $E_i=0$。对应欧拉示性数 $[E]=\sum(-1)^i[E_i]$，便得到自然满射
\[
\mathcal{E}(X)\longrightarrow K(X).
\]
给出 $K(X)$ 的 $\mathcal{E}(X)$ 的等价关系，显然由同构与基本对象的加法生成。这里，基本对象是指对某个 $j$ 满足
\[
P_j=P_{j+1},\qquad P_i=0\quad(i\ne j,j+1)
\]
的形如 $\sum P_i$ 的对象。同样，对 $G$ 空间 $X$，也可以把 $K_G(X)$ 表示成 $\mathcal{E}_G(X)$ 的商。其中，$\mathcal{E}_G(X)$ 的元素是分次 $G$ 向量丛。

% [原文 p.137]
这里设 $E\in\mathcal{E}(X)$ 为分次向量丛。此时 $E^{\otimes k}$ 也是分次向量丛，其次数按惯例定义为 $k$ 个因子次数之和。$E^{\otimes k}$ 上 $S_k$ 的作用不只有因子的置换，还要一并考虑适当的符号。即，二项 $e_p\otimes e_q$（$e_p\in E_p$，$e_q\in E_q$）的互换带上符号 $(-1)^{pq}$。此时，$E^{\otimes k}$ 的欧拉示性数 $[E^{\otimes k}]$ 是 $K_{S_k}(X)$ 的元素。

% [Tutor Note: 交换 $p$ 次与 $q$ 次因子时乘 $(-1)^{pq}$ 是 Koszul 符号规则，正是为了让置换作用成为复形的自同构：否则 $d^{\otimes k}$ 与重排不对易，下面证明中“$Q$ 作为复形非轮状”就不成立。这一步是易错点。]

\begin{proposition*}{B.2.2}
  元素 $[E^{\otimes k}]\in K_{S_k}(X)$ 只依赖于 $E$ 的同构类 $[E]\in K(X)$。因此确定如下算子
  \[
  \otimes k\colon K(X)\longrightarrow K_{S_k}(X)=K(X)\otimes R(S_k).
  \]
\end{proposition*}

\begin{proof}
  设 $P$ 为 $\mathcal{E}(X)$ 的基本对象，则须证明
  \[
  [(E\oplus P)^{\otimes k}]=[E^{\otimes k}]\in K_{S_k}(X).
  \]
  然而，有 $S_k$ 表示的分解
  \[
  (E\oplus P)^{\otimes k}\cong E^{\otimes k}\oplus Q,
  \]
  故只需证 $K_{S_k}(X)$ 中 $[Q]=0$。为此，把 $E$ 看作所有映射均为零的向量丛复形，把 $P$ 看作具有恒等映射 $P_j\to P_{j+1}$ 的复形。此时 $(E\oplus P)^{\otimes k}$ 是向量丛的复形，$S_k$ 作为复形自同构映射之群作用于其上（因为符号正是为此而选的）。对 $E^{\otimes k}$ 与 $Q$ 也有同样的事实。这里 $Q$ 以 $P$ 为因子，故 $Q$ 确实非轮状（完全）。因此，由本节开头所注意的事项，正如所求，$K_{S_k}(X)$ 中 $[Q]=0$。
\end{proof}

\begin{remark}{}
  令 $\pi(E)=\Hom_{S_k}(V_\pi,E^{\otimes k})$，并用 $S_k$ 的作用把 $E^{\otimes k}$ 分解为
  \[
  E^{\otimes k}\cong\sum_{\pi}V_\pi\otimes\pi(E),
  \]

% [原文 p.138]
如此分解时，命题 B.2.2 断言 $E\mapsto\pi(E)$ 诱导算子
\[
\pi\colon K(X)\longrightarrow K(X).
\]
\end{remark}

以 $\operatorname{Op}(K)$ 表示函子 $K$ 到自身的所有自然变换之集。亦即，元素 $T\in\operatorname{Op}(K)$ 对每个 $X$ 确定映射
\[
T(X)\colon K(X)\longrightarrow K(X),
\]
并对拉回是自然的。$\operatorname{Op}(K)$ 的加法与乘法分别由值的加法与乘法定义。即，对 $a\in K(X)$，令
\begin{gather*}
(T+S)(X)(a)=T(X)(a)+S(X)(a),\\
TS(X)(a)=T(X)(a)\cdot S(X)(a).
\end{gather*}

把算子
\[
\otimes k\colon K(X)\longrightarrow K(X)\otimes R(S_k)
\]
与同态 $\varphi\colon R(S_k)\to\Z$ 复合，便得到自然映射
\[
T_\varphi\colon K(X)\longrightarrow K(X).
\]
这定义了群同态
\[
j_k\colon R_*(S_k)\longrightarrow\operatorname{Op}(K).
\]
把它作加法扩张，便得到环同态
\[
j\colon R_*\longrightarrow\operatorname{Op}(K).
\]
这样便达到了说明对称群如何定义 K 理论算子这一目的。此外，这个环 $R_*$ 的结构已在 B.1 节完全决定。在结束本节之前，先考察几个特殊的算子，
% [原文 p.139]
把本文中的定义与 Grothendieck（cf. [5, §12]）以及 Adams（[2]）给出的定义联系起来。

为避免复杂的式子，通常省略记号 $j$，而把 $R_*$ 的元素直接看作算子。实际上，证明 $j$ 是单射同态并不困难，故可以把 $R_*$ 看作 $\operatorname{Op}(K)$ 的子环。不过，本文并不需要这一事实。

现在，B.1 节所述的 $\sigma^k$、$\lambda^k$、$\psi^k$、$\theta^p$ 都可以看作 K 理论的算子。由定义可知，若 $E$ 为向量丛，则 $\lambda^k[E]$ 是 $E$ 的 $k$ 次外幂的等价类，而 $\sigma^k(E)$ 是 $E$ 的 $k$ 次对称幂的等价类。一般地，$K(X)$ 的任一元素都可表为 $[E_0]-[E_1]$（$E_0,E_1$ 为向量丛）。把 $(E_0\oplus E_1)^{\otimes k}$ 看作 $S_k$ 复形，取出对称部分与反对称部分，便知
\begin{equation}
\sigma^k\bigl([E_0]-[E_1]\bigr)=\sum_{j=0}^{k}(-1)^j\sigma^{k-j}[E_0]\lambda^j[E_1] \tag{1}
\end{equation}
\begin{equation}
\lambda^k\bigl([E_0]-[E_1]\bigr)=\sum_{j=0}^{k}(-1)^j\lambda^{k-j}[E_0]\sigma^j[E_1] \tag{2}
\end{equation}
成立。以 $u$ 为不定元，形式地令 $\lambda_u=\sum\lambda^ku^k$，$\sigma_u=\sum\sigma^ku^k$，并在式 (1) 中令 $E_0=E_1$，则得
\begin{equation}
\sigma_u[E_1]\lambda_{-u}[E_1]=1. \tag{3}
\end{equation}
这个等式当然也可以用命题 B.1.2 的同构映射，从生成 $e_k$ 与 $h_k$ 的函数之间的对应关系推出。这里，由式 (2) 得
\begin{align*}
\lambda_u\bigl([E_0]-[E_1]\bigr)&=\lambda_u[E_0]\sigma_{-u}[E_1]\\
&=\lambda_u[E_0]\lambda_u[E_1]^{-1}.
\end{align*}
其中最后一等号用了 (3)。本来 Grothendieck 就是用这个式子把 $\lambda^k$ 从向量丛扩张到 K 群的。因此，在本文中把 $\lambda^k$ 定义为算子的做法与 Grothendieck 的定义一致。用分次张量积把算子从向量丛扩张到 K 群的方法，是 Grothendieck
% [原文 p.140]
用于外幂的方法的推广\footnote{这一事实 Grothendieck 无疑知道。}。

Adams 用 Newton 多项式，令
\[
\psi^k=Q_k(\lambda^1,\dots,\lambda^k),
\]
把 Grothendieck 的 $\lambda^k$ 定义为算子 $\psi^k$。因此由系 B.1.8，本文的定义与 Adams 的定义一致。$\psi^k$ 的重要性质是加法性。因此，下面说明怎样从本文的定义直接证明它。

% [Tutor Note: $Q_k$ 是 $k$ 次 Newton 多项式，它把初等对称函数 $e_1,\dots,e_k$ 与幂和 $p_k=\sum t_i^k$ 互化；Adams 正是用它由外幂 $\lambda^k$ 反解出 $\psi^k$。所以“$\psi^k$ 由 $\lambda^1,\dots,\lambda^k$ 表示”不是定义选择，而是 Newton 恒等式的直接推论。]

\begin{proposition*}{B.2.3}
  设 $E,F$ 为向量丛，则
  \[
  \psi^k\bigl([E]\pm[F]\bigr)=\psi^k[E]\pm\psi^k[F]
  \]
  成立。
\end{proposition*}

\begin{proof}
  取 $D_0=E$、$D_1=F$ 的分次向量丛 $D$，考虑 $D^{\otimes k}$。按与命题 B.1.1 相同的论证，可得
  \[
  [D]^{\otimes k}=\sum_{j=0}^{k}(-1)^j\operatorname{ind}_j[E^{\otimes k-j}\otimes F^{\otimes j}]\in K(X)\otimes R(S_k).
  \]
  其中 $\operatorname{ind}_j\colon K(X)\otimes R(S_{k-j}\times S_j)\to K(X)\otimes R(S_k)$ 由诱导表示给出。这里 $E^{\otimes k-j}$ 按标准作用为 $S_{k-j}$ 向量丛，而 $S_j$ 按置换与次数的符号作用于 $F^{\otimes j}$。为得到 $\psi^k[D]$，必须对长 $k$ 的循环置换求值 $R(S_k)$。与命题 B.1.1 相同，除 $j=0,k$ 之外的各项均为零。长 $k$ 的循环置换的符号为 $(-1)^{k-1}$，故
  \begin{align*}
  \psi^k\bigl([E]-[F]\bigr)&=\psi^k[E]+(-1)^k(-1)^{k-1}\psi^k[F]\\
  &=\psi^k[E]-\psi^k[F]
  \end{align*}
  成立。对 $[E]+[F]$ 也可同样处理，但更简单。
\end{proof}

\para{乘法性}

% [原文 p.141]
\[
\psi^k[E\otimes F]=\psi^k[E]\psi^k[F]
\]
由同型
\[
(E\otimes F)^{\otimes k}\cong E^{\otimes k}\otimes F^{\otimes k}
\]
与迹的乘法性立得。

这里与系 B.1.4 相同，把基本要素 $\Delta_{n,k}$ 展开为
\[
\Delta_{n,k}=\sum a_i\otimes b_i.
\]
其中 $a_i\in\operatorname{Sym}_k[t_1,\dots,t_n]$ 是基，从而 $b_i\in R(S_k)$ 是对偶基（假定 $n\ge k$）。此时，相应地，对任意 $x\in K(X)$，得到 $x^{\otimes k}$ 的展开
\[
x^{\otimes k}=\sum_i\alpha_i(x)\otimes b_i\in K(X)\otimes R(S_k).
\]
其中 $\alpha_i=(\Delta')^{-1}a_i\in R_*$。这由 $\Delta'$ 的定义与 $R_*$ 对 $K(X)$ 的作用方式立刻可知。

设 $a_i$ 为由基本对称式构成的单项式，则 $\alpha_i$ 是由对应的外幂 $\lambda^i$ 构成的单项式。因此由命题 B.1.9 可知下面的命题\footnote{现在既然把 B.1 节的 $\lambda^i$ 与外幂等同看待，便采用通常的记法：不写 $\Lambda^i(M)$ 而写 $\lambda^i(M)$，相应地不写 $\Lambda_{-i}(M)$ 而写 $\lambda_{-i}(M)$。}。

\begin{proposition*}{B.2.4}
  对任意 $x\in K(X)$，有
  \[
  x^{\otimes k}=(-1)^{k-1}\lambda^k(x)\otimes\lambda_{-1}(M)+\text{合成项}.
  \]
  其中“合成”是指含有至少两个 $\lambda^i(x)$ 的乘积，$M$ 是命题 B.1.9 所定的 $S_k$ 的 $(k-1)$ 维表示。
\end{proposition*}

这里考虑循环群 $\Z/k$。把 $K(X)\otimes R(\Z/k)$ 中 $x^{\otimes k}$ 的像记作 $P^k(x)$，称为 $k$ 次循环幂（cyclic power）。在 $k=p$（素数）的特殊情形，由 (B.1.11) 导出下面的命题。

% [原文 p.142]
\begin{proposition*}{B.2.5}
  设 $p$ 为素数，$x\in K(X)$。此时 $p$ 次循环幂 $P^p(x)$ 由式
  \[
  P^p(x)=\psi^p(x)\otimes1+\theta^p(x)\otimes N\in K(X)\otimes R(\Z/p)
  \]
  给出。其中 $N$ 是 $\Z/p$ 的正则表示。
\end{proposition*}

这里，由同构映射
\[
\Delta'\colon R_*\longrightarrow\varprojlim_{n}\operatorname{Sym}[t_1,\dots,t_n],
\]
$\psi^p$ 与 $\theta^p$ 分别对应多项式 $\sum t_i^p$ 与 $\dfrac{(\sum t_i)^p-\sum t_i^p}{p}$。因此，它们由式
\[
\psi^p=(\psi^1)^p-p\theta^p
\]
相联系，且对任意 $x\in K(X)$，
\[
\psi^p(x)=x^p-p\theta^p(x)
\]
成立。把它代入命题 B.2.5 的式子，便得到
\begin{equation}
P^p(x)=x^p\otimes1+\theta^p(x)\otimes(N-p). \tag{B.2.6}
\end{equation}
这比命题 B.2.5 的写法更好。因为以 $I(\Z/p)$ 为增广理想时，它对应分解
\[
R(\Z/p)=\Z\oplus I(\Z/p).
\]
因此，
\[
\theta^p(x)\otimes(N-p)\in K(X)\otimes I(\Z/p)
\]
表示 $p$ 次循环幂 $P^p(x)$ 与「通常」的 $p$ 次幂 $x^p\otimes1$ 之差。

命题 B.2.5 对向量丛 $V$ 给出 $\psi^p[V]$ 的简明而几何的描述。设 $T$ 为 $V^{\otimes p}$ 的因子的循环置换，$V_j$ 为对应于特征值 $\exp(2\pi\sqrt{-1}j/p)$ 的 $T$ 的特征空间，则

% [原文 p.143]
\begin{equation}
\psi^p[V]=[V_0]-[V_1] \tag{B.2.7}
\end{equation}
成立。实际上，由命题 B.2.5 可知
\[
[V_0]=\psi^p[V]+\theta^p[V],\qquad [V_j]=\theta^p[V]\quad(j=1,\dots,p-1).
\]

\section{外部张量幂}
\takeaway{本节把算子 $\otimes k$ 相对化、外部化：用 $G$ 复形定义相对群 $K_G(X,Y)$，再把 $k$ 重张量幂搬到直积空间上，得到外部张量幂算子 $\boxtimes k\colon K(X,Y)\to K_{S_k}(X,Y)^{k}$。}

为进一步考察算子 $\otimes k$ 的性质，需要把它“相对化”与“外部化”。

首先，设 $X$ 为 $G$ 空间，$Y$ 为 $X$ 的部分 $G$ 空间，考虑相对 K 群 $K_G(X,Y)$。与通常 K 群的情形相同，可以把 $K_G(X,Y)$ 看作集合 $\mathcal{E}_G(X,Y)$ 按等价关系取商。$\mathcal{E}_G(X,Y)$ 的对象 $E$ 是 $X$ 上、在 $Y$ 上非轮状的向量丛 $G$ 复形。亦即，$E$ 由 $G$ 向量丛 $E_i$（除有限个 $i$ 外 $E_i=0$）与同 $G$ 的作用可交换的同态映射序列
\[
\cdots\longrightarrow E_i\xrightarrow{\ d\ }E_{i+1}\xrightarrow{\ d\ }\cdots
\]
组成，且 $d^2=0$，并使得该序列在 $Y$ 的每点上完全。基本对象 $P$ 是指 $P_i=0$（$i\ne j,j+1$）、$P_j=P_{j+1}$ 且 $d\colon P_j\to P_{j+1}$ 为恒等映射者。赋予 $\mathcal{E}_G(X,Y)$ 的等价关系由同构与基本对象的加法（直和）生成。此时，对 $E\in\mathcal{E}_G(X,Y)$，把其等价类记作 $[E]\in K_G(X,Y)$。细节请参看 [4]。无群作用情形的类似结果，即把 $K(X,Y)$ 定义为 $\mathcal{E}(X,Y)$ 之商，请参看 [7, part II]。

% [Tutor Note: 这一套与 §2.6 的 $\mathcal{C}_n$、$\mathcal{L}_n$ 是同一台机器的等变版本，区别只在：这里允许任意长的复形（不固定长度 $n$），且所有映射须与 $G$ 的作用可交换。基本对象 $P$ 的加法就是 §2.6 中的“基本列”。]

其次，考虑外部张量幂。设 $E$ 为 $X$ 上的向量丛，定义 $E^{\boxtimes k}$ 为直积 $X^k$（$k$ 个 $X$ 的直积）上的向量丛，使点 $(x_1,x_2,\dots,x_k)$ 上的纤维为 $E_{x_1}\otimes E_{x_2}\otimes\cdots\otimes E_{x_k}$。$E^{\boxtimes k}$
% [原文 p.144]
是 $S_k$ 空间 $X^k$ 上的 $S_k$ 向量丛。

这里，对称群 $S_k$ 通过置换直积因子作用于 $X^k$。

设
\[
\delta\colon X\longrightarrow X^k
\]
为对角映射，则显然有自然的 $S_k$ 同构映射
\begin{equation}
\delta^*(E^{\boxtimes k})\cong E^{\otimes k}. \tag{B.3.1}
\end{equation}

若 $E$ 是 $X$ 上向量丛的复形，同样可以定义 $E^{\boxtimes k}$，它是 $X^k$ 上向量丛的复形。此外，$X^k$ 如前所述是 $S_k$ 空间，$E^{\boxtimes k}$ 是向量丛的 $S_k$ 复形。若 $E$ 在 $Y\subset X$ 上非轮状，则在由至少一个 $i$ 满足 $x_i\in Y$ 的点 $(x_1,x_2,\dots,x_k)$ 构成的 $X^k$ 的子空间上，$E^{\boxtimes k}$ 非轮状。把这个子空间记作 $X^{k-1}Y$，并把对 $(X^k,X^{k-1}Y)$ 记作 $(X,Y)^k$。这样便定义了算子
\[
\boxtimes k\colon\mathcal{E}(X,Y)\longrightarrow\mathcal{E}_{S_k}(X,Y)^k.
\]

命题 B.2.2 的证明可直接推广到这一情形，于是下面的命题成立。

\begin{proposition*}{B.3.2}
  算子 $E\mapsto E^{\boxtimes k}$ 诱导算子
  \[
  \boxtimes k\colon K(X,Y)\longrightarrow K_{S_k}(X,Y)^k.
  \]
\end{proposition*}

\begin{corollary*}{B.3.3}
  若 $x$ 含于 $K(X)\to K(Y)$ 的核中，则 $x^{\boxtimes k}$ 含于
  \[
  K_{S_k}(X^k)\longrightarrow K_{S_k}(X^{k-1}Y)
  \]
  的核中。
\end{corollary*}

\begin{proof}
  这由命题 B.3.2 与算子 $\boxtimes k$ 的自然性立得。
\end{proof}

由 (B.3.1)，得到下面的交换图。

% CHUNK appB2 DONE
